% De Haruspicum Responsis --- headnote
% de-haruspicum-responsis, late May / early summer 56 BC, Rome

\textit{On the Responses of the Haruspices, delivered in the
Senate at Rome late in May or early summer of 56 BC. The
occasion is a haruspical response on a strange noise
``with a rumbling'' (\textit{strepitus cum fremitu})
heard in the agro Latiniensi, the open country south of
Rome --- which the Etruscan diviners had explained as a
warning that ``games had been less diligently performed
and polluted'' (referring to the Megalesia of April 56,
recently put on by Clodius as curule aedile, when he
loosed slaves into the Senate's seats), that ``sacred
and religious places are being held profane'' (which
Clodius read as a reference to Cicero's recovered
Palatine house, which he claimed was still consecrated
ground), and various other faults; and that the
remedy was concord among the optimates, lest discord
``bring slaughter and danger to the chief men, and the
state pass to one rule.'' Clodius had spoken in the
contio that morning, twisting the response into a
prosecution of Cicero's house. Cicero answers in the
Senate the same day, picking apart Clodius's reading
phrase by phrase and showing every clause to be a
prosecution not of him but of Clodius.}

\textit{The structure is the most architecturally
controlled of Cicero's mature speeches. \textsection
1--17 is the apology for the previous day, when Cicero
had cut off Clodius mid-question with a threat of trial
and chased him out of the Curia (the famous tableau
\textsection 2: Clodius collapses on the Curia threshold
on seeing Cn. Lentulus Marcellinus the consul, ``out of
recollection of his Gabinius and longing for his Piso'').
\textsection 18--19 is the great patriotic argument from
religion: the forefathers as masters of religion, the
sense in which the Romans surpass the Spaniards in number,
the Gauls in strength, the Phoenicians in cleverness, the
Greeks in arts, the Italians and Latins in inborn sense
--- by piety alone. \textsection 20--29 takes the
haruspical response head by head: polluted games are the
Megalesia (the Magna Mater rumble in the agro Latiniensi
read against the Phrygian rite at Pessinus that Clodius
sold to the Galatian Brogitarus, against Deiotarus's
recovery of the priesthood, and the long-standing dignity
of the rite). \textsection 30--33 returns to the house:
mine has been freed by every judgement; yours --- where
you killed Q. Seius and where the chapel still stands ---
is the polluted one. \textsection 34--39 (envoys killed,
faith and oath neglected, ancient and secret rites)
returns again and again to Clodius: the Alexandrians,
Plator of Orestis bled by the physician at Thessalonica,
the bribed jurors of the Bona Dea trial, and the great
prosecution of the Bona Dea outrage itself --- closing
on the famous answer to the question what punishment the
gods inflict: not pain or wound of body but frenzy and
madness.}

\textit{\textsection 40--55 turns to the second half of
the response, ``the discord of the optimates,'' and uses
it as the frame for the great catalogue of popularis
turncoats: Tiberius and Gaius Gracchus, Saturninus,
Sulpicius --- each had a grave (if not just) cause; only
Clodius was made popularis ``out of the saffron-coloured
robe, the headband, the women's slippers, the breast-band,
the lyre.'' Sections 47--52 expose the men who let
Clodius continue as a tribune-attack-dog on Pompey: the
threat to loose Caesar's army on the Senate, the dagger
in the temple of Castor at the Pompey siege, the late
turn (when Clodius now praises Pompey, having sold himself
in turn to every camp). The closing \textsection 53--55
makes the haruspical word \textit{principes} explicit:
the gods foretell danger to the chief men --- Pompey,
Crassus, Caesar, the senatorial \textit{principes} ---
not to Clodius, who is no \textit{princeps}.}

\textit{The peroration \textsection 56--63 covers the
remaining heads of the response: the inferior and the
rejected, the warning against change in the standing
of the state, and (\textsection 60) the famous catalogue
of the symptoms of the late Republic at this very year:
the empty treasury, the publicans not enjoying their
contracts, the authority of the chief men fallen, the
consensus of the orders broken, the courts perished, the
votes held by a few. The closing prescription is
concord: not human counsels but divine. The thirty days
between this speech and the conference of Luca, when
Caesar reset the terms of the triumvirate behind Cicero's
back, give the speech its tragic edge --- it is the last
of the post-exile speeches in the optimate-Pompeian-Senate
register, and the last in which Cicero can speak as if
\textit{principes} are a stable category in Roman public
life.}
